\documentclass[12pt]{article}
\usepackage[utf8]{inputenc}
\usepackage[T1]{fontenc}
\usepackage{hyperref}
\title{Nisi Language Resources for K-12 Curriculum Development}
\author{Compiled by Claude Sonnet 5.5 (claude-sonnet-5-5) for Andr\'as Kornai}
\date{}
\begin{document}
\maketitle

\section{Introduction}
Nisi (ISO 639-3 \texttt{yso}, Glottocode nisi1238) is a Sino-Tibetan (Loloish/Ngwi) language of Yunnan, China. Glottolog places it in Southeastern Ngwi, as an unclassified member of that group. It is not known to be taught in schools. The web search found no dedicated Nisi materials. Nisi may be related to the Nisu varieties described by other codes (nsd, yiv, nos, nsv, nsf), but I could not confirm that. No resource below should be treated as a Nisi (yso) resource unless it is checked against that code.

\section{Primary Educational Resources}
No primers, textbooks or school programmes for Nisi (yso) were found.

\section{Linguistic and Reference Materials}
\begin{itemize}
\item Glottolog entry for Nisi (China), nisi1238, with its classification: \url{https://glottolog.org/resource/languoid/id/nisi1238.bigmap.html}
\item Wikipedia, Nisu language. This covers the related Southern Yi (Nisu) varieties, not specifically yso: \url{https://en.wikipedia.org/wiki/Nisu_language}
\item Wikipedia, Nyisu language: \url{https://en.wikipedia.org/wiki/Nyisu_language}
\item Lama, dissertation on the subgrouping of the Nisoic languages (University of Texas at Arlington), a comparative study of 34 Nisoic languages: \url{https://uta-ir.tdl.org/uta-ir/bitstream/handle/10106/11161/Lama_uta_2502D_11591.pdf%3Fsequence=1&amp;isAllowed=y}
\end{itemize}

\section{Religious and Community Resources}
Joshua Project lists Bible translation as needed for the related Nisu varieties. These are not confirmed to be yso.
\begin{itemize}
\item Southwestern Nisu (nsv): \url{https://joshuaproject.net/languages/nsv}
\item Northwestern Nisu (nsf): \url{https://joshuaproject.net/languages/nsf}
\item Eastern Nisu (nos): \url{https://joshuaproject.net/index.php/languages/nos}
\end{itemize}

\section{Digital and Audio Resources}
\begin{itemize}
\item Wikitongues page for Eastern Nisu (nos), a related variety: \url{https://wikitongues.org/languages/nos/}
\end{itemize}
No audio or corpus resources specific to yso were found.

\section{Curriculum Development Context}
There is currently no basis for a native-language curriculum from published materials. The first steps would be to fix the identity and dialect of yso, then document its lexicon and orthography with community members. Related Nisu varieties may offer comparative material only after it is verified.
\end{document}